% Source: https://artofproblemsolving.com/community/c3h1396926_geometry
% Stable identity: aops-topic-1396926-q01
The mark on a wheel touches the ground exactly when the carriage has gone a whole number of turns of that wheel. So we need the least positive distance $d$ with
\[
d=\dfrac{7}{2}p=\dfrac{11}{3}q
\]
for positive integers $p$, $q$. Multiplying by $6$: $21p=22q$. Since $\gcd(21,22)=1$, $22\mid p$; the least choice is $p=22$, $q=21$, giving $d=\dfrac{7}{2}\cdot22=\boxed{77}$ metres.

Why the numerators: a common multiple of $\dfrac{7}{2}$ and $\dfrac{11}{3}$ must be a multiple of $7$ and of $11$ after clearing the denominators. In general, for fractions in lowest terms, the least common multiple is
\[
\operatorname{lcm}\left(\dfrac{a}{b},\dfrac{c}{e}\right)=\dfrac{\operatorname{lcm}(a,c)}{\gcd(b,e)},
\]
and here $\dfrac{\operatorname{lcm}(7,11)}{\gcd(2,3)}=\dfrac{77}{1}=77$. (The answer is the product $7\cdot11$ only because $7$ and $11$ are coprime and the denominators are coprime.)
